% !TeX root = ../tesis.tex

\section{Método}
\label{sec:metodo}
%-----------------------------------------------------------------------

La presente investigación adopta un enfoque \textbf{cuantitativo} fundamentado
en la teoría de grafos y el modelado de redes complejas de transporte. Este
enfoque permite representar formalmente la red multimodal de la \zmvm{} como un
grafo dirigido y ponderado $\grafo = (\nodos, \aristas, \peso{i}{j})$, sobre el
cual se calculan indicadores de eficiencia, vulnerabilidad y conectividad de
manera objetiva y reproducible \citep{allauca2023, owais2020}.

La elección de este método responde a tres condiciones estructurales de la
investigación: (1)~la disponibilidad de datos abiertos estandarizados en formato
\gtfs{} para la totalidad de los sistemas de transporte masivo de la \cdmx{};
(2)~la existencia de motores computacionales de autoría propia ---
\textit{Apimetro} \citep{cabrera2025apimetro} y \textit{VFTModel}
\citep{cabrera2026VFTModel}--- que permiten procesar dichos datos a escala
metropolitana; y (3)~la naturaleza cuantificable de las preguntas de
investigación, centradas en la reducción de tiempos de traslado, la distribución
de centralidad y la robustez geométrica de la red.

A diferencia de los enfoques tradicionales de modelación basados en matrices
origen-destino, el análisis mediante grafos no requiere calibración de encuestas
de movilidad, lo que lo hace especialmente adecuado para la evaluación
estructural de redes existentes y la simulación de escenarios topológicos
alternativos \citep{fortin2016innovative}.

Las herramientas y fuentes de datos que articulan el proceso metodológico son:

\begin{itemize}
\item \textbf{Feeds \gtfs{}} de la red de Movilidad Integrada de la \cdmx{} y
    los sistemas del Estado de México (Mexibús, Mexicable), como fuente primaria
    de datos estructurados de la red de transporte.
\item \textit{Apimetro} \citep{cabrera2025apimetro}: motor de adquisición,
normalización y construcción del grafo multimodal georreferenciado a partir de
    datos \gtfs{}.
\item \textit{VFTModel} \citep{cabrera2026VFTModel}: motor analítico para el
    cálculo de indicadores de eficiencia urbana ($DI$, $T$), centralidad de
    intermediación ($B(v)$) y robustez geométrica ($\Delta E$) sobre el grafo
    construido.
\item \textit{Transport-gis} \citep{cabrera2026transportgis}: motor de
    visualización cartográfica para la generación de mapas temáticos
    georreferenciados a partir de los indicadores calculados por
    \textit{VFTModel}, empleando capas \gis{} en formato GeoPackage,
    QGIS y Tableau como plataformas de representación territorial.
\item \textbf{Python y Go (Golang)}: lenguajes de programación empleados en el
    procesamiento de datos, análisis espacial y desarrollo de los motores
    computacionales.
    \item \textbf{Sistemas de Información Geográfica (\gis{})}: QGIS y GeoPandas
para la visualización, validación espacial y análisis territorial de resultados.
\end{itemize}

%-----------------------------------------------------------------------
\section{Metodología}
\label{sec:metodologia}
%-----------------------------------------------------------------------

La investigación se estructura en cuatro fases secuenciales e interdependientes,
cada una orientada a cubrir uno o más de los objetivos específicos planteados.
El hilo conductor del proceso es la construcción progresiva del análisis: de la
adquisición y modelado de los datos hacia el diagnóstico estructural de la red,
la simulación de escenarios y, finalmente, la formulación de lineamientos de
política pública en materia de movilidad \citep{allauca2023}.

\subsection{Fase 1: Adquisición y construcción del grafo multimodal}
\label{subsec:fase1}

Esta fase tiene una ponderación del 25\% sobre el total de la investigación y
cubre el \textbf{Objetivo Específico~1}.

\begin{description}
\item[Objetivo de la Fase:] Construir el grafo multimodal georreferenciado de
    la \zmvm{} a partir del procesamiento sistemático de los feeds \gtfs{} de la
    totalidad de los sistemas de transporte masivo de la región.

    \item[Actividades a Desarrollar:]~
    \begin{itemize}
        \item \textit{Adquisición y normalización de datos \gtfs{}}: Descarga y
        limpieza de los feeds oficiales mediante el motor \textit{Apimetro}
\citep{cabrera2025apimetro}. Este proceso incluye la validación de la
        integridad estructural de los archivos (\texttt{stops.txt},
\texttt{routes.txt}, \texttt{trips.txt}, \texttt{stop\_times.txt}) conforme
        al estándar \gtfs{} \citep{fortin2016innovative}.
\item \textit{Construcción del grafo dirigido y ponderado}: Transformación
de los datos normalizados en un grafo $\grafo = (\nodos, \aristas)$, donde
        cada nodo $v \in \nodos$ representa una estación o parada y cada arista
        $e \in \aristas$ un segmento de ruta ponderado por tiempo de viaje
        $\peso{i}{j}$.
        \item \textit{Georreferenciación y validación espacial}: Asociación de
coordenadas geográficas a cada nodo y verificación de la topología resultante
        mediante herramientas \gis{}.
    \end{itemize}
\end{description}

\subsection{Fase 2: Análisis de la topología actual con \textit{VFTModel}}
\label{subsec:fase2}

Esta fase tiene una ponderación del 30\% y cubre los
\textbf{Objetivos Específicos~2 y~3}.

\begin{description}
    \item[Objetivo de la Fase:] Caracterizar cuantitativamente la estructura y
    eficiencia de la red de transporte actual mediante el cálculo de indicadores
    de teoría de redes sobre el grafo construido en la Fase~1.

    \item[Actividades a Desarrollar:]~
    \begin{itemize}
\item \textit{Cálculo del Índice de Ruta Directa ($DI$) y Tiempo de Viaje
        Promedio ($T$)}: Aplicación del motor \textit{VFTModel}
        \citep{cabrera2026VFTModel} para obtener los indicadores base de
eficiencia en la configuración topológica actual. El $DI$ mide la desviación
de los trayectos reales respecto a la distancia euclidiana óptima; valores
elevados indican ineficiencia geométrica en la red \citep{owais2020}.
        \item \textit{Cálculo de la Centralidad de Intermediación ($B(v)$)}:
Identificación de los nodos con mayor flujo de rutas mínimas, localizando
los puntos neurálgicos cuya falla comprometería la conectividad global del
        sistema \citep{batista2015, loterovélez2014}.
\item \textit{Evaluación de la Robustez Geométrica ($\Delta E$)}: Cálculo
del impacto sistémico ante la eliminación simulada de nodos o aristas de alta
centralidad, cuantificando la degradación de la eficiencia global de la red
        \citep{loterovélez2014}.
    \end{itemize}
\end{description}

\subsection{Fase 3: Diagnóstico de cobertura y simulación de escenarios}
\label{subsec:fase3}

Esta fase tiene una ponderación del 35\% y cubre los
\textbf{Objetivos Específicos~4 y~5}.

\begin{description}
    \item[Objetivo de la Fase:] A partir del diagnóstico estructural obtenido en
    la Fase~2, identificar las principales brechas de cobertura y eficiencia del
    sistema actual, y evaluar el impacto topológico de la incorporación de un
    corredor anillar en el Tercer Contorno metropolitano.

    \item[Actividades a Desarrollar:]~
    \begin{itemize}
\item \textit{Diagnóstico de cobertura y eficiencia}: Análisis espacial de
los indicadores calculados para identificar zonas con déficit de conectividad
        lateral, tiempos de traslado elevados y dependencia de nodos de alta
        intermediación \citep{mishra2012, ugur2025}.
\item \textit{Diseño del escenario propuesto}: Modelado de corredores anillares
en los cuatro cuadrantes del Anillo Periférico como un conjunto de nuevos
nodos $\nodos'$ y aristas $\aristas'$ incorporados al grafo original, con
pesos $\peso{i}{j}$ estimados a partir de la velocidad comercial de sistemas
análogos de transporte masivo.
\item \textit{Simulación comparativa}: Ejecución de \textit{VFTModel} sobre
        el grafo ampliado para calcular la variación neta en $DI$, $T$ y $B(v)$,
cuantificando la reducción de tiempos de traslado y la redistribución de la
        centralidad de intermediación \citep{cabrera2026VFTModel}.
    \end{itemize}
\end{description}

\subsection{Fase 4: Síntesis e interpretación de resultados}
\label{subsec:fase4}

Esta fase tiene una ponderación del 10\% y complementa el
\textbf{Objetivo Específico~5}.

\begin{description}
\item[Objetivo de la Fase:] Interpretar los resultados comparativos de la
simulación hipotética para evaluar la capacidad analítica de la metodología
desarrollada y establecer las condiciones de replicabilidad del proceso en
otros contextos metropolitanos con datos \gtfs{} disponibles.

    \item[Actividades a Desarrollar:]~
    \begin{itemize}
        \item \textit{Lectura comparativa de indicadores}: Análisis de la
        variación neta en $DI$, $T$ y $B(v)$ entre la red actual y la red
        ampliada, valorando la magnitud del impacto topológico del escenario
        anillar hipotético.
        \item \textit{Documentación de la metodología}: Registro de los
        parámetros, supuestos y limitaciones del proceso analítico, con el
        fin de garantizar su reproducibilidad en contextos similares.
    \end{itemize}
\end{description}
